% =============================================================================
% 第7章：讨论——启示、局限与未来工作
% =============================================================================

\section{讨论}
\label{sec:discussion}

第\ref{sec:results}节和第\ref{sec:diagnosis}节分别从"现象"和"机制"两个层面完成了注意力塌缩的诊断：PCA空间的嵌入同构是塌缩的根因，Focal Loss的梯度聚焦将塌缩放大了$\times$3.3倍，Sparsemax通过稀疏投影将注意力多样性恢复了$\times$7.5倍。本节将这一发现置于更广阔语境中讨论。

\subsection{对表格深度学习的启示}

\textbf{（1）自注意力不是表格数据的通用解。}自注意力的有效性高度依赖于输入表征的语义丰富性。在NLP中，每个token对应一个具有独立含义的词——注意力通过比较词语义身份来分配权重。在表格数据中，当特征具有明确的业务语义（如"交易金额"、"商户类别"）时，Feature Tokenizer可以将这些语义映射为嵌入方向。但当特征经过PCA匿名化——丧失了语义身份——注意力便丧失了其运作的前提条件。本文的诊断证据表明，这一失效不是"性能比树模型差几个点"的温和退化，而是注意力机制在功能上（均值池化）和表征上（多头同构）的系统性塌缩。

\textbf{（2）PCA的数据处理决策应与模型选择联动。}在实际的机器学习管线中，数据预处理（如PCA匿名化）和模型选择通常是独立决策——数据团队做PCA以保护隐私或降维，建模团队选择模型架构。本文的诊断结果提示：如果业务需求决定了数据必须经过PCA匿名化，那么使用自注意力模型需要额外的验证步骤（至少应检查注意力的分布特性），或者直接选择不依赖跨特征语义关联的架构。

\textbf{（3）PCA匿名化场景比表面看起来更普遍。}虽然本文的分析基于信用卡欺诈检测这一特定数据集，但PCA匿名化作为一种"特征语义真空"的创造机制，其适用范围远超此单一场景。在联邦学习中，出于隐私保护目的的特征变换（如差分隐私梯度扰动、安全聚合后的特征空间扭曲）可能产生与PCA类似的"语义均一化"效应。在多传感器融合场景中，来自同质传感器阵列的测量值（如多个加速度计、多个光谱波段）也具有PCA匿名化特征的结构特性——数值分布相似、个体"身份"仅由通道编号区分。在差分隐私保护的表格数据发布中，特征可能经过随机投影或PCA降维以混淆个体信息。这些场景的共同特点是：特征丧失了可独立识别的语义身份，而这正是本文诊断的注意力塌缩的结构性前提。因此，本文的发现和诊断框架的适用范围可能远广于"信用卡欺诈检测+PCA"这一特定组合。

\textbf{（4）标准化也是语义抹平的来源——来自Amount锚点验证的负结果。}本文设置了一项补充分析：使用数据集中的Amount特征（唯一未经PCA处理的特征，保留了"交易金额"的业务语义）作为"语义锚点"，检验其嵌入方向是否比V1-V28更独立。结果显示，Amount的嵌入范数仅排名20/29，其与V特征的余弦相似度（0.0214）与V特征之间的余弦相似度（0.0163）无显著差异。这意味着\textbf{在本数据集中，即使是具有原始业务语义的特征，在StandardScaler标准化后也表现出与PCA特征相似的方向同构性}——标准化将不同量纲、不同分布的业务特征强制统一为均值0、标准差1的分布，可能在此过程中压缩了特征的"身份信息"。由于这一观察仅基于单个特征（Amount）且差异微小，尚不能作为强结论，但提示了一个值得进一步检验的假设：塌缩的结构性前提可能不仅限于PCA匿名化，任何使特征数学属性高度同构的处理流程（标准化、PCA、白化）都可能削弱自注意力的前提条件。

\textbf{（3）诊断优于比较。}本文的一个隐含方法论立场是：对于深度学习机制的理解，系统\textbf{诊断}一个模型的失效模式比\textbf{比较}多个模型的性能数字提供更多可迁移的知识。FTT+Softmax的AUPRC=0.9126这一数字本身没有揭示"发生了什么"；是注意力分布、嵌入SVD谱和多头多样性的诊断分析，揭示了塌缩的本质和Sparsemax修复的边界。这一立场对学术研究具有操作性的指导意义：如果一篇论文的核心叙事是"方法A优于方法B"，其结论的适用范围受限于所测试的数据集；如果核心叙事是"方法A在条件X下以机制Y失效"，其结论是可迁移的——任何满足条件X的设置都预期观察到模式Y。

\subsection{实践的选型指南}

基于本文的诊断发现，下表总结了在表格数据上选择模型时的条件考量：

\begin{table}[htbp]
\centering
\caption{基于特征语义丰富度的模型选型建议。此表仅为本文诊断发现的定性推论，未经系统性实验验证。}
\label{tab:selection_guide}
\begin{tabular}{p{3.5cm}p{5.5cm}p{5cm}}
\toprule
\textbf{特征类型} & \textbf{推荐方案} & \textbf{理由} \\
\midrule
PCA匿名化/均一化特征 & 避免自注意力；优先MLP或树模型 & 语义真空→嵌入同构→注意力塌缩风险 \\
有业务语义的原始特征 & 自注意力可用，但建议先做注意力分布检查 & 语义差异为嵌入层提供区分性信号 \\
混合（部分匿名+部分有语义） & 分离处理：有语义特征用注意力，匿名特征用MLP & 避免匿名特征的噪声干扰注意力聚集 \\
\bottomrule
\end{tabular}
\end{table}

\subsection{局限性与未来工作}

本文存在以下主要局限：

\begin{enumerate}
    \item \textbf{单数据集}：所有分析基于一个PCA匿名化的信用卡欺诈数据集。结论在以下两类数据上的泛化性有待验证：（a）非匿名化的表格数据（特征具有独立业务语义）——预期自注意力不会发生同类塌缩；（b）其他类型的特征同构数据，如传感器阵列、光谱数据、基因组数据。
    \item \textbf{单一种子}：所有结果基于seed=42的单次训练，无法区分真实效应与随机初始化差异。未来工作应以多种子（$\geq$5）训练，报告AUPRC和注意力分布的均值与标准差。
    \item \textbf{BCE+Sparsemax缺失}：本文发现了BCE下注意力塌缩减轻的现象（std 0.00408 vs 0.00125），但未检验Sparsemax在BCE下的额外增益。这一组合是"不平衡处理—注意力多样性"权衡的最直接检验——BCE+Sparsemax是否能同时保留BCE的注意力自由度和Sparsemax的稀疏选择——是未来工作的最高优先级。
    \item \textbf{Sparsemax的未解耦效应}：Sparsemax同时改变了注意力输出的稀疏性和熵——两种效应在本文中未被分离。温度缩放（temperature scaling）+Softmax可能以不同机制实现类似的熵调节效果。此外，近期文献中提出的替代方案（如可调节稀疏度的$\alpha$-Entmax、ICML 2024的MultiMax、以及表格专用的BiSHop）可能提供比固定约束Sparsemax更灵活的稀疏-多峰平衡，值得在后续研究中系统比较。
    \item \textbf{诊断框架的外部验证}：第\ref{sec:diagnosis}节提出的三阶段诊断框架（注意力分布$\to$嵌入SVD$\to$因果归因）在本文中仅在一个数据集上演示，应被视为"提出的诊断模板"而非"已验证的通用框架"。该框架在其他数据集和模型上的操作性和预测有效性需要独立验证。
\end{enumerate}

未来工作方向包括：（1）在非匿名化表格数据集上复制本文的诊断流程，检验"语义丰富度$\to$嵌入多样性$\to$注意力区分度"的理论链；（2）训练BCE+Sparsemax模型以直接检验"不平衡处理—注意力多样性"权衡；（3）通过温度缩放消融和$\alpha$-Entmax实验分离Sparsemax的稀疏化与熵调节效应；（4）探索嵌入层预训练策略（如对比学习），在类似PCA空间中恢复特征嵌入的方向多样性——从根本上解决本文诊断的"语义真空"问题。
